% ---------------------------------------------------------------------------
% Standalone section: a run of the sign-flipped BKNS algorithm that gets stuck.
% Uses only macros already in preamble.tex (\items, \agents, \cost, \subsidy,
% \optsubsidy, \set, \abs, \Z, \Mset). Drop into sections/ and \input it.
% ---------------------------------------------------------------------------

\section{The goods algorithm, sign-flipped, gets stuck}
\label{sec:whynotbkns}

\subsection{The instance}

Three agents $\agents = \set{1,2,3}$ and three chores $\items = \set{a,b,c}$.
Each cost function depends only on the size of the bundle:
\[
  \cost_1(S) \;=\; \max\bigl(0,\ \abs S - 1\bigr),
  \qquad
  \cost_2(S) \;=\; \cost_3(S) \;=\; \min\bigl(\abs S,\,2\bigr).
\]
Every marginal is $0$ or $1$ and $\cost_i(\emptyset) = 0$, so the instance is
negative dichotomous.

\begin{center}
\begin{tabular}{c|cccc}
$\abs S$ & $0$ & $1$ & $2$ & $3$ \\
\hline
$\cost_1(S)$ & $0$ & $0$ & $1$ & $2$ \\
$\cost_2(S) = \cost_3(S)$ & $0$ & $1$ & $2$ & $2$ \\
\end{tabular}
\end{center}

\subsection{The algorithm, sign-flipped}

Algorithm~1 of~\cite{BKNS22} allocates one item per iteration and maintains an
envy-free solution with subsidies in $\set{0,1}^n$. Its EXTEND branch requires
the incoming item to have marginal $+1$ for the recipient. No chore has a
positive marginal, so under the sign flip that branch never applies and every
iteration goes to FINDSINK, which assigns the item to an agent of
\emph{maximum} subsidy, i.e.\ an agent in $\Mset(\subsidy)$.

\subsection{The run}

\begin{center}
\begin{tabular}{llcc}
step & allocation & $\optsubsidy$ & $\Mset(\optsubsidy)$ \\
\hline
start & $(\emptyset,\ \emptyset,\ \emptyset)$ & $(0,0,0)$ & $\set{1,2,3}$ \\
$a$ to agent $1$ & $(\set{a},\ \emptyset,\ \emptyset)$ & $(0,0,0)$ & $\set{1,2,3}$ \\
$b$ to agent $1$ & $(\set{a,b},\ \emptyset,\ \emptyset)$ & $(1,0,0)$ & $\set{1}$ \\
\end{tabular}
\end{center}

Both steps are legal: at each one the recipient lies in $\Mset(\optsubsidy)$
and the resulting allocation is envy-freeable with subsidies in $\set{0,1}^3$.

\subsection{The last chore}

At this point $\Mset(\optsubsidy) = \set 1$, so FINDSINK must give $c$ to agent
$1$. It cannot. Nor can any other agent, and nor does reassigning the bundles
help:

\begin{center}
\begin{tabular}{lll}
$c$ added to & resulting bundles & least $\optsubsidy$ over all reassignments \\
\hline
agent $1$'s bundle & $\set{a,b,c},\ \emptyset,\ \emptyset$ & $(2,0,0)$ \\
agent $2$'s bundle & $\set{a,b},\ \set{c},\ \emptyset$ & $(2,1,0)$ \\
agent $3$'s bundle & $\set{a,b},\ \emptyset,\ \set{c}$ & $(2,0,1)$ \\
\end{tabular}
\end{center}

Every option requires a subsidy of $2$. The algorithm is stuck.

\subsection{The instance itself is fine}

The allocation $(\set a,\ \set b,\ \set c)$ gives every agent cost $1$ and is
envy-free with $\subsidy = \mathbf 0$. What fails is the one-item-at-a-time
construction: reaching it requires taking back the two chores already placed on
agent $1$.
